% GENERATED FILE. Edit canonical lessons, then run npm run generate:books.
% canonical-source-hash: fnv1a64:4662b2f8b17fc3f2
% canonical-lessons: ZH-C241-zhuyuan, ZH-C241-bingfang, ZH-C241-shoushu, ZH-C241-chuyuan, ZH-C241-huifu-recover, ZH-R241-review-236-241

\chapter{In Hospital}
\label{ch:zh-in-hospital}
\hlchaptermodality{241}

\begin{chapteropening}
\textbf{By the end of this chapter:} \emph{I can tell how someone went into hospital, the ward, the operation, coming home and getting better.}

Complete the last lesson of chapter 241: I can tell how someone went into hospital, the ward, the operation, coming home and getting better.
\end{chapteropening}

\section[zhùyuàn]{zhùyuàn — to be in hospital}

\label{lesson:ZH-C241-zhuyuan}



\begin{quote}
Before the new one: say the Mandarin for to throw up, then the Mandarin for to have an upset stomach.
\end{quote}

\subsection*{You'll want to know: zhùyuàn}
\emph{zhùyuàn} — \textquotedblleft{}to be in hospital\textquotedblright{}.

\begin{quote}
\emph{shàng ge yuè wǒ bàba zhùyuàn le} — Last month my dad went into hospital.
\end{quote}

Its characters come later; for now it is a word for the ear.

\begin{cousinweb}
\emph{zhùyuàn} is built from \emph{zhù} 'to stay' and \emph{yuàn} 'a walled courtyard', both inherited from Middle Chinese, \emph{drjuH} and \emph{hjwenH} (Guangyun, in Baxter's transcription), which glosses the first 'to stop' and the second 'a wall round a courtyard'.
\end{cousinweb}

\begin{grammarlens}[title={What you've built}]
One more word for talking about the past.
\end{grammarlens}

\subsection*{Your turn}
\begin{itemize}
\raggedright
  \item \emph{Say it:} \emph{zhùyuàn}
  \item \emph{Say it:} \emph{zhùyuàn}, once more
  \item \emph{Say it:} \emph{shàng ge yuè wǒ bàba zhùyuàn le} — Last month my dad went into hospital.
  \item \emph{Recall:} say the Mandarin for to throw up, then the Mandarin for to have an upset stomach, then say \emph{zhùyuàn} again
\end{itemize}

\begin{tcolorbox}[breakable,colback=teal!4,colframe=teal!35!black,title={Before you move on}]
What does zhùyuàn mean? (\textquotedblleft{}To be in hospital\textquotedblright{}.) Say it once more.
\end{tcolorbox}
\section[\texorpdfstring{\zh{病房} (bìngfáng)}{病房 (bìngfáng)}]{\zh{病房} (bìngfáng) — hospital ward}

\label{lesson:ZH-C241-bingfang}



\begin{quote}
Before the new one: say the Mandarin for to have an upset stomach, then the Mandarin for to be in hospital.
\end{quote}

\subsection*{You'll want to know: \zh{病房}}
\textbf{\zh{病房}} — \emph{bìngfáng} — \textquotedblleft{}hospital ward\textquotedblright{}.

\begin{quote}
\textbf{\zh{医生现在在病房里}} \emph{yīshēng xiànzài zài bìngfáng lǐ} — The doctor's in the ward now.
\end{quote}

\begin{cousinweb}
\textbf{\zh{病房}} (\emph{bìngfáng}) is built from \zh{病} 'illness' and \zh{房} 'room', both inherited from Middle Chinese, \emph{bjaengH} and \emph{bjang} (Guangyun, in Baxter's transcription), which glosses the second 'a chamber'.
\end{cousinweb}

\begin{grammarlens}[title={What you've built}]
One more word for talking about the past.
\end{grammarlens}

\subsection*{Your turn}
\begin{itemize}
\raggedright
  \item \emph{Say it:} \emph{bìngfáng}
  \item \emph{Say it:} \emph{bìngfáng}, once more
  \item \emph{Say it:} \emph{yīshēng xiànzài zài bìngfáng lǐ} — The doctor's in the ward now.
  \item \emph{Recall:} say the Mandarin for to have an upset stomach, then the Mandarin for to be in hospital, then say \emph{bìngfáng} again
\end{itemize}

\begin{tcolorbox}[breakable,colback=teal!4,colframe=teal!35!black,title={Before you move on}]
What does \zh{病房} mean? (\textquotedblleft{}Hospital ward\textquotedblright{}.) Say it once more.
\end{tcolorbox}
\section[shǒushù]{shǒushù — surgical operation}

\label{lesson:ZH-C241-shoushu}



\begin{quote}
Before the new one: say the Mandarin for to be in hospital, then the Mandarin for a hospital ward.
\end{quote}

\subsection*{You'll want to know: shǒushù}
\emph{shǒushù} — \textquotedblleft{}surgical operation\textquotedblright{}.

\begin{quote}
\emph{tā zuòle yí cì shǒushù} — He had an operation.
\end{quote}

Its characters come later; for now it is a word for the ear.

\begin{cousinweb}
\emph{shǒushù} is built from \zh{手} 'hand' and \emph{shù} 'skill', both inherited from Middle Chinese, \emph{syuwX} and \emph{zywit} (Guangyun, in Baxter's transcription), which glosses the second 'skill'.
\end{cousinweb}

\begin{grammarlens}[title={What you've built}]
One more word for talking about the past.
\end{grammarlens}

\subsection*{Your turn}
\begin{itemize}
\raggedright
  \item \emph{Say it:} \emph{shǒushù}
  \item \emph{Say it:} \emph{shǒushù}, once more
  \item \emph{Say it:} \emph{tā zuòle yí cì shǒushù} — He had an operation.
  \item \emph{Recall:} say the Mandarin for to be in hospital, then the Mandarin for a hospital ward, then say \emph{shǒushù} again
\end{itemize}

\begin{tcolorbox}[breakable,colback=teal!4,colframe=teal!35!black,title={Before you move on}]
What does shǒushù mean? (\textquotedblleft{}Surgical operation\textquotedblright{}.) Say it once more.
\end{tcolorbox}
\section[chūyuàn]{chūyuàn — to leave hospital}

\label{lesson:ZH-C241-chuyuan}



\begin{quote}
Before the new one: say the Mandarin for a hospital ward, then the Mandarin for a surgical operation.
\end{quote}

\subsection*{You'll want to know: chūyuàn}
\emph{chūyuàn} — \textquotedblleft{}to leave hospital\textquotedblright{}.

\begin{quote}
\emph{tā xià xīngqī kěyǐ chūyuàn} — He can leave hospital next week.
\end{quote}

Its characters come later; for now it is a word for the ear.

\begin{cousinweb}
\emph{chūyuàn} is built from \zh{出} 'to go out' and \emph{yuàn} 'a walled courtyard', both inherited from Middle Chinese, \emph{tsyhwit} and \emph{hjwenH} (Guangyun, in Baxter's transcription).
\end{cousinweb}

\begin{grammarlens}[title={What you've built}]
One more word for talking about the past.
\end{grammarlens}

\subsection*{Your turn}
\begin{itemize}
\raggedright
  \item \emph{Say it:} \emph{chūyuàn}
  \item \emph{Say it:} \emph{chūyuàn}, once more
  \item \emph{Say it:} \emph{tā xià xīngqī kěyǐ chūyuàn} — He can leave hospital next week.
  \item \emph{Recall:} say the Mandarin for a hospital ward, then the Mandarin for a surgical operation, then say \emph{chūyuàn} again
\end{itemize}

\begin{tcolorbox}[breakable,colback=teal!4,colframe=teal!35!black,title={Before you move on}]
What does chūyuàn mean? (\textquotedblleft{}To leave hospital\textquotedblright{}.) Say it once more.
\end{tcolorbox}
\section[huīfù]{huīfù — to recover; to get back to normal}

\label{lesson:ZH-C241-huifu-recover}



\begin{quote}
Before the new one: say the Mandarin for a surgical operation, then the Mandarin for to leave hospital.
\end{quote}

\subsection*{You'll want to know: huīfù}
\emph{huīfù} — \textquotedblleft{}to recover; to get back to normal\textquotedblright{}.

\begin{quote}
\emph{tā de shēntǐ huīfù de hěn kuài} — He's recovering very quickly.
\end{quote}

Its characters come later; for now it is a word for the ear.

\begin{cousinweb}
\emph{huīfù} is built from \emph{huī} 'vast' and \emph{fù} 'to return', both inherited from Middle Chinese, \emph{khwoj} and \emph{bjuwk} (Guangyun, in Baxter's transcription), which glosses the first 'great' and the second 'to return'. The h of huī is not the regular outcome of that form.
\end{cousinweb}

\begin{grammarlens}[title={What you've built}]
One more word for talking about the past.
\end{grammarlens}

\subsection*{Your turn}
\begin{itemize}
\raggedright
  \item \emph{Say it:} \emph{huīfù}
  \item \emph{Say it:} \emph{huīfù}, once more
  \item \emph{Say it:} \emph{tā de shēntǐ huīfù de hěn kuài} — He's recovering very quickly.
  \item \emph{Recall:} say the Mandarin for a surgical operation, then the Mandarin for to leave hospital, then say \emph{huīfù} again
\end{itemize}

\begin{tcolorbox}[breakable,colback=teal!4,colframe=teal!35!black,title={Before you move on}]
What does huīfù mean? (\textquotedblleft{}To recover; to get back to normal\textquotedblright{}.) Say it once more.
\end{tcolorbox}
\section[(dialogue)]{Looking back: from Something Odd to In Hospital}

\label{lesson:ZH-R241-review-236-241}



\begin{quote}
Say the Mandarin for to turn up and to find out.
\end{quote}

\subsection*{Your turn}
\begin{itemize}
\raggedright
  \item \emph{Say it:} to turn up, to find out, to be scared, luckily, nearly
  \item \emph{Say it:} a plant, a peanut, a leaf, to sprout, a watermelon
  \item \emph{Say it:} the sky, the moon, dark clouds, the open sea, a rainbow
  \item \emph{Say it:} healthy, to give up smoking, to drink alcohol, to breathe, a medical check-up
  \item \emph{Say it:} to get sick, to become inflamed, swollen, to throw up, to have an upset stomach
  \item \emph{Say it:} to be in hospital, a hospital ward, a surgical operation, to leave hospital, to recover
  \item \emph{Read it:} \zh{现} — each one aloud, then a word that holds it
\end{itemize}

\begin{tcolorbox}[breakable,colback=teal!4,colframe=teal!35!black,title={Before you move on}]
To turn up? (\textbf{\zh{出现}}, \emph{chūxiàn}.) To find out? (\emph{fāxiàn}.) To be scared? (\emph{hàipà}.) Luckily? (\emph{xìnghǎo}.) Nearly? (\emph{chàdiǎnr}.) A plant? (\emph{zhíwù}.) A peanut? (\textbf{\zh{花生}}, \emph{huāshēng}.) A leaf? (\emph{yèzi}.) To sprout? (\emph{fāyá}.) A watermelon? (\emph{xīguā}.) The sky? (\emph{tiānkōng}.) The moon? (\emph{yuèliang}.) Dark clouds? (\emph{wūyún}.) The open sea? (\emph{dàhǎi}.) A rainbow? (\emph{cǎihóng}.) Healthy? (\emph{jiànkāng}.) To give up smoking? (\emph{jièyān}.) To drink alcohol? (\emph{hējiǔ}.) To breathe? (\emph{hūxī}.) A medical check-up? (\emph{tǐjiǎn}.) To get sick? (\textbf{\zh{生病}}, \emph{shēngbìng}.) To become inflamed? (\emph{fāyán}.) Swollen? (\emph{zhǒng}.) To throw up? (\emph{tù}.) To have an upset stomach? (\emph{lā dùzi}.) To be in hospital? (\emph{zhùyuàn}.) A hospital ward? (\textbf{\zh{病房}}, \emph{bìngfáng}.) A surgical operation? (\emph{shǒushù}.) To leave hospital? (\emph{chūyuàn}.) To recover? (\emph{huīfù}.)
\end{tcolorbox}
